\section{Domain-Adaptive Pre-Training}

\label{sec:pretraining}
% ══════════════════════════════════════════════════════════════════════════════

Standard ImageNet pre-training provides suboptimal features for
fine-grained ecological recognition due to domain mismatch (controlled
studio photos vs.\ field conditions) and category mismatch (1{,}000
general categories vs.\ thousands of visually similar species).

\subsection{Pre-Training Data}

We assemble \textbf{EcoVision-28M}, a dataset of 28 million ecological
images from:

\begin{itemize}
  \item \textbf{iNaturalist 2021}: 10 million images across 10{,}000
    species \citep{van2021benchmarking}.
  \item \textbf{GBIF media}: 8 million specimen and field images.
  \item \textbf{Museum collections}: 5 million digitized specimen
    photographs from 14 natural history museums.
  \item \textbf{Camera trap archives}: 3 million images from Wildlife
    Insights and LILA BC \citep{beery2021iwildcam}.
  \item \textbf{Citizen science platforms}: 2 million images from
    eBird, BugGuide, and Reef Life Survey.
\end{itemize}

\subsection{Self-Supervised Pre-Training}

We pre-train ViT-L/14 using DINOv2 \citep{oquab2024dinov2} on
EcoVision-28M for 300{,}000 iterations with batch size 1{,}024. The
self-supervised objective learns representations that capture
fine-grained visual structure without requiring species labels.

\subsection{Supervised Fine-Tuning}

After self-supervised pre-training, we fine-tune on the labeled subset
(iNaturalist 2021, 10M images, 10{,}000 species) using a hierarchical
classification head with losses at family, genus, and species level:
\begin{equation}
  \Loss_{\text{pretrain}} = \Loss_{\text{species}} + 0.3
  \Loss_{\text{genus}} + 0.1 \Loss_{\text{family}}
\end{equation}


% ══════════════════════════════════════════════════════════════════════════════
